\section{Iterated cycle removal without an excess parameter}
\label{sec:four-iterative-packing}

This argument uses the all-terminal forest estimate in clean regions.
It needs no lower bound on a positive excess above one half.
The marked vertex set must be connected or empty.

Fix $H=H_4$ from Proposition~\ref{prop:four-fixed-ratio-shore-rank}.
Choose an integer $D_*\ge3$ so large that Corollary~\ref{cor:four-fixed-ratio-terminal-forests} gives forest probability at most $3/5$ in every affine terminal class of every normalized pole with
\[
 h\ge H,\qquad 2\le k\le4D,\qquad \operatorname{girth}\ge D\ge D_*.
\]
These are absolute constants.

\begin{lemma}[Two clean regions]
\label{lem:four-two-clean-regions}
Let $G$ be connected and cubic, let $P$ be connected or empty, and let $S$ be disjoint from $P$.
Suppose $G-(P\cup S)$ has girth at least $D\ge D_*$.
If $R_1,R_2$ are disjoint connected vertex sets avoiding $P\cup S$, each of rank at least $H$ and boundary at most $4D$, then
\[
 \Pr(\xi_G-(P\cup S)\text{ is a forest})\le\frac9{25}.
\]
\end{lemma}

\begin{proof}
Prune leaves from each induced graph $G[R_i]$.
The resulting connected core has the same rank and no larger boundary.
Indeed, removing a leaf in a cubic ambient graph reduces the boundary by one.
Each core is a normalized pole: its vertices have internal degrees two or three, and each degree-two vertex has one boundary edge.

Condition on every edge outside the two core interiors.
For every feasible conditioning, the internal states are independent and uniform in their respective affine parity classes.
This follows because the fiber is the product of two affine translates of their cycle spaces.
Each boundary pattern has even size and is supported on the degree-two ports.
If a core has at least two ports, its conditional forest probability is at most $3/5$ by the choice of $D_*$.
With zero or one port, the only feasible boundary pattern is empty, and its forest probability is $2^{-h}\le3/5$.
The global forest event requires both internal states to be forests.
Their conditional independence gives the bound $(3/5)^2$.
\end{proof}

\begin{proposition}[Packing after a previous deletion]
\label{prop:four-packing-after-deletion}
Under the preceding notation, suppose
\[
 \Pr(\xi_G-(P\cup S)\text{ is a forest})>\frac12,
 \qquad s=|S|.
\]
Let $2\le d\le2D$.
Every family of $q$ disjoint connected cyclic regions in $G-(P\cup S)$, each with boundary in $G$ at most $d$, satisfies
\[
 q<(6H+15)d(s+1)2^d.
\]
\end{proposition}

\begin{proof}
The case $q=0$ is immediate.
Contract the regions and the connected components outside them, discarding loops, to obtain a connected quotient $Q$.
As in Proposition~\ref{prop:four-excess-free-cycle-packing}, the forest probability is at most
\[
 \mathbb E\prod_v(1-I_v/2),
\]
where $I_v$ indicates a zero star at a marked quotient vertex.
No local star space is zero.
A maximal direct-sum family of local spaces has a vertex set $A$ with
\[
 1\le t=|A|<2^d.
\]
If it includes all marked vertices, the result follows.
Otherwise put $B=|\delta_QA|\le dt$.
Lemma~\ref{lem:four-terminal-bipartition} gives at most $2B$ exceptional remaining vertices and at most $2B$ ordered classes with identical terminal bipartitions.

Consider one class $v_1,\ldots,v_m$ in its tree order.
Every consecutive block of $H+2$ class vertices gives a middle component after deleting its first and last vertices from $Q-A$.
That component contains its $H$ interior marked vertices.
It has no neighbor in $A$, since the first and last vertices give the same actual terminal bipartition.
Its boundary in $Q$ therefore ends only at these two vertices and has size at most $2d$.
The middle components from disjoint consecutive blocks are disjoint.
Their lifts to $G$ are connected and have rank at least $H$.

The connected set $P$ lies within one component outside all marked regions, hence within one quotient vertex.
Thus at most one of these disjoint lifts meets $P$.
At most $s$ further lifts meet $S$.
If $m\ge(H+2)(s+3)$, there are at least $s+3$ middle components and at least two clean lifts.
Each has boundary at most $2d\le4D$.
Lemma~\ref{lem:four-two-clean-regions} would give forest probability at most $9/25$, a contradiction.
Consequently each class has fewer than $(H+2)(s+3)$ members.

It follows that
\begin{align*}
 q&\le t+2B+2B(H+2)(s+3)\\
  &\le t+(6H+14)(s+1)B\\
  &\le(6H+15)d(s+1)t
   <(6H+15)d(s+1)2^d.
\end{align*}
\end{proof}

\begin{theorem}[An excess-free girth-raising bound]
\label{thm:four-excess-free-girth-raising}
There is an absolute constant $K$ with the following property.
Let $G$ be connected and cubic and let $P$ be connected or empty.
Assume
\[
 \Pr(\xi_G-P\text{ is a forest})>\frac12,
 \qquad \operatorname{girth}(G-P)\ge D_*.
\]
For every integer $D\ge D_*$, there is $S\subseteq V(G)\setminus P$ such that $G-(P\cup S)$ has girth greater than $D$ and
\[
 |S|+1\le 2^{\,2D+K\log^2(D+2)}.
\]
Moreover, $S$ is the union of a vertex-disjoint family of cycles of length at most $D$, maximal among such cycles outside $P$.
For two target scales $D_0<D_1$, they can be chosen nested with the respective logarithmic bounds
\[
 2D_0+K\log^2(D_0+2),\qquad
 2D_1+2D_0+2K\log^2(D_1+2).
\]
\end{theorem}

\begin{proof}
Deletion can only increase the probability of the forest event.
Starting with $S=\varnothing$, repeatedly raise the scale from $a$ to $b\le2a$ by taking a maximal disjoint family of cycles of length at most $b$ in $G-(P\cup S)$ and adding their vertices to $S$.
The remaining graph has girth greater than $b$.
If the family is nonempty, its cycles have boundary at most $b$ in the cubic ambient graph.
Proposition~\ref{prop:four-packing-after-deletion}, with $d=b$, gives
\[
 s_{\mathrm{new}}+1
 \le (6H+16)b^2 2^b(s+1).
\]
The same inequality holds when the family is empty.

Choose the scales backwards from $D$, using $\lceil D/2^j\rceil$ until reaching a value between $D_*$ and $2D_*$.
Read them in increasing order, starting from the available girth $D_*$.
If $D=D_*$, perform one deletion at that same scale.
Each forward step at most doubles the previous scale.
The sum of the new scales is at most $2D+O(\log(D+2))$, and there are $O(\log(D+2))$ steps.
Taking logarithms of the recurrence gives
\[
 \log(s+1)
 \le2D+O\!\left(\sum_b\log(b+2)\right)
 \le2D+K\log^2(D+2).
\]
The fixed initial scale is absorbed into $K$.
Each new family avoids all earlier deletions, so the combined cycles are vertex-disjoint.
Their lengths are at most $D$, and the final girth assertion is exactly their maximality outside $P$.
For two targets, first carry out the construction through $D_0$.
Retain those deletions, then use the backwards-planned scales for $D_1$ down to the first value between $D_0$ and $2D_0$.
The second sum of new scales is at most $2D_1+O(\log(D_1+2))$.
Adding the two logarithmic bounds gives the stated nested version.
\end{proof}

\paragraph{Scope.}
The theorem starts above the fixed girth threshold $D_*$.
It does not bound the number of short cycles below that threshold.
Its two-scale consequence is $|S_0|=r^{o(1)}$ at $D_0=\log r/\log\log r$ and $|S_1|\le r^{2\alpha+o(1)}$ at $D_1=\alpha\log r$, for any fixed $0<\alpha<1/2$.

\paragraph{Independent checks.}
The mathematical proof has passed independent line audits by the root, path-stability, and random-host agents.
The standard-library checker \texttt{four-high-connectivity-iterative-packing-check.py} additionally verifies actual terminal partitions, 636 middle intervals, 114 disjoint segment pairs, 4482 core prunings, and 4163 exact conditional affine fibers.
Its results are recorded in \texttt{four-high-connectivity-iterative-packing-checks.json}.
These finite checks test the structural assertions and conditional product law; they do not replace the high-girth theorem.
